% =====================================================================
\phantomsection
\addcontentsline{toc}{part}{The Main Language of the Course}
\BookChapter[Python]{Python: The Main Language of the Course}
\label{ch:python}

\begin{researchquestion}
What do you need to know about Python to write, test and trust a small engineering program?
\end{researchquestion}
% =====================================================================

\begin{factsheet}{Python}
First appeared & 1991 (\py~3.0 in 2008; a new version, 3.x, appears every October)\\
Designed by & Guido van Rossum, Centrum Wiskunde \& Informatica, Netherlands\\
Paradigms & procedural, object-oriented, functional; array style through NumPy\\
Typing & dynamic, strong; optional type hints\\
Execution & interpreted (CPython bytecode interpreter); alternatives include PyPy (JIT)\\
File extension & \file{.py}; notebooks \file{.ipynb}\\
Cost & free and open source (PSF licence)\\
Used for & data analysis, scientific computing, machine learning, automation,
  scripting of engineering software, web back ends, teaching\\
Famous software & NumPy, SciPy, pandas, Matplotlib, scikit-learn, PyTorch, TensorFlow,
  Jupyter, ASE, pymatgen, FEniCS, Blender scripting, Abaqus scripting\\
\end{factsheet}

\py is the main language of ENG 101. This chapter is a compact but complete
introduction that follows the order of the course. It is longer than the other language
chapters because it is the reference point for all of them: when a later chapter says
``unlike \py'', this is the \py it means.

\LearningOutcomes{Write scripts that use variables, decisions, loops, collections and functions; read and write data files; handle errors with exceptions; organise related data and functions in a class; compute with NumPy arrays instead of loops; plot results with Matplotlib; call SciPy and pandas for standard numerical and data tasks; and debug and test your own code.}

\section{Why Python?}

\begin{history}
Guido van Rossum started \py as a hobby project over the Christmas holidays of 1989,
while working at the Centrum Wiskunde \& Informatica in Amsterdam. He wanted a language
that was easy to read, easy to extend with C code, and fun to use. The name comes from the
British comedy group Monty Python, not from the snake. \py~2.0 (2000) added many features
that made it popular; \py~3.0 (2008) cleaned up the language at the cost of breaking
compatibility, and the transition took more than a decade. \py~2 reached its official
end of life in January 2020. Only \py~3 is used today.
\end{history}

\py's success in engineering and science rests on a few things.
\begin{itemize}
\item \textbf{Readability.} \py code looks almost like pseudocode. The language's design
  principles, which you can print with \code{import this}, include ``Readability
  counts'' and ``There should be one --- and preferably only one --- obvious way to do
  it''.
\item \textbf{Batteries included.} The standard library contains modules for mathematics,
  file formats, dates, networking, testing and much more.
\item \textbf{The scientific ecosystem.} NumPy, SciPy, Matplotlib, pandas and Jupyter give
  \py most of what MATLAB offers, for free. On top of them sit thousands of specialised
  packages for every branch of engineering.
\item \textbf{A glue language.} \py can call libraries written in C, C++ and Fortran
  easily, and many commercial engineering programs (Abaqus, ANSYS, Aspen, COMSOL, ArcGIS,
  Blender, FreeCAD, ParaView) can be controlled with \py scripts.
\item \textbf{Machine learning.} The major deep-learning frameworks (PyTorch, TensorFlow,
  JAX) are driven from \py.
\end{itemize}
Its main weakness is speed: pure-\py loops are slow (Section~\ref{sec:numpy-speed}).
Section~\ref{sec:python-strengths} discusses this and other limits.

\section{Getting started}

\subsection*{Install a distribution that includes the scientific libraries}

There are several good ways to get \py:
\begin{itemize}
\item \textbf{Anaconda} or its smaller sibling \textbf{Miniforge} install \py together
  with NumPy, SciPy, Matplotlib, pandas and Jupyter in one step, and include the
  \code{conda} package manager. This is the easiest option for engineering students.
\item \textbf{python.org} offers the official installer for Windows and macOS. You then
  install packages with \code{pip}.
\item On Linux, \py is usually already installed; use the system package manager or
  \code{uv}/\code{pip} in a virtual environment.
\item \textbf{Online notebooks} such as Google Colab and JupyterLite run \py in a web
  browser with nothing to install.
\end{itemize}

\subsection*{Run Python interactively, as a script, or in a notebook}

\paragraph{Interactive mode.} Type \code{python} (or \code{python3}) in a terminal. The
prompt \code{>>>} appears, and every line you type is executed immediately:
\begin{lstlisting}[language=Python]
>>> 2 + 3
5
>>> import math
>>> math.sqrt(80 * 9.81 / 0.25)
56.02856414093976
\end{lstlisting}
This is perfect for quick calculations and for trying out an idea.

\paragraph{Scripts.} For anything longer, write the code in a text file with the extension
\file{.py} and run the whole file:
\codefile{Python}{python/hello.py}{a first script}
\begin{lstlisting}[style=shell]
$ python3 hello.py
\end{lstlisting}
\outfile{python/hello.out}

\paragraph{Notebooks.} A \term{Jupyter notebook} (\file{.ipynb}) mixes code cells,
their output (numbers, tables, plots) and formatted text with equations in one document.
Notebooks are ideal for exploring data, documenting an analysis and teaching; many ENG
101 case studies are notebooks. Their weakness is that cells can be run in any order, so
a notebook can end up in a state that nobody can reproduce. Before you share a notebook,
choose ``Restart kernel and run all cells'' and check that it still works.

\subsection*{Use an editor that can run and debug code}

An \term{IDE} (integrated development environment) combines an editor, a way to run code,
and a \term{debugger} that lets you pause a program and inspect its variables. Popular
choices for \py are \textbf{VS Code} (free, works with every language in these notes),
\textbf{Spyder} (MATLAB-like, included in Anaconda), \textbf{PyCharm} and
\textbf{JupyterLab}. Use whichever your lab sessions use; the code is the same.

\section{Values, variables and types}

\subsection*{A variable is a name attached to a value}

A \term{variable} is a name that refers to a value. You create it by \term{assignment}
with a single equals sign:
\begin{lstlisting}[language=Python]
mass = 80.0              # kg
g = 9.81                 # m/s^2
weight = mass * g        # N
print(weight)            # 784.8000000000001
\end{lstlisting}
The \code{=} sign does not mean mathematical equality. It means ``compute the value on
the right, then attach the name on the left to it''. That is why \code{n = n + 1} makes
sense in a program even though it is false in algebra.

Variable names may contain letters, digits and underscores, must not start with a digit,
and are case sensitive (\code{Mass} and \code{mass} are different). By convention \py
uses \code{lower\_case\_with\_underscores}. Choose names that say what the value means:
\code{flow\_rate} is better than \code{q}, and much better than \code{x2}. Some words are
reserved by the language and cannot be used as names: \code{if}, \code{for}, \code{def},
\code{class}, \code{lambda}, \code{True}, \code{None}, and a few others.

\subsection*{Six basic types cover most needs}

\begin{center}
\small
\begin{tabular}{@{}llll@{}}
\toprule
\textbf{Type} & \textbf{Meaning} & \textbf{Examples} & \textbf{Notes}\\
\midrule
\code{int} & integer & \code{42}, \code{-7}, \code{1\_000\_000} & unlimited size\\
\code{float} & real number & \code{9.81}, \code{6.022e23}, \code{-0.5} & IEEE double\\
\code{complex} & complex number & \code{3+4j} & \code{j} is $\sqrt{-1}$\\
\code{bool} & truth value & \code{True}, \code{False} & \\
\code{str} & text & \code{"steel"}, \code{\sq{}Al\sq{}}, \code{"""..."""} & Unicode\\
\code{NoneType} & ``no value'' & \code{None} & \\
\bottomrule
\end{tabular}
\end{center}
The function \code{type()} tells you the type of a value, and \code{int()},
\code{float()}, \code{str()} convert between types:
\begin{lstlisting}[language=Python]
>>> type(9.81)
<class 'float'>
>>> int("42") + 1
43
>>> float("2.5e3")
2500.0
>>> str(3.0) + " m"
'3.0 m'
\end{lstlisting}

\subsection*{Division always gives a float, and \texttt{**} is the power}

\begin{center}
\small
\begin{tabular}{@{}lll@{}}
\toprule
\textbf{Operator} & \textbf{Meaning} & \textbf{Example}\\
\midrule
\code{+ - *} & add, subtract, multiply & \code{2*3 - 1} $\to$ \code{5}\\
\code{/} & true division (always a float) & \code{7/2} $\to$ \code{3.5}\\
\code{//} & floor division & \code{7//2} $\to$ \code{3}, \code{-7//2} $\to$ \code{-4}\\
\code{\%} & remainder (modulo) & \code{7 \% 2} $\to$ \code{1}\\
\code{**} & power & \code{2**10} $\to$ \code{1024}, \code{2**0.5} $\to$ \code{1.414\ldots}\\
\bottomrule
\end{tabular}
\end{center}
Operators follow the usual mathematical precedence: \code{**} first, then \code{* / // \%},
then \code{+ -}. When in doubt, use parentheses. Note that \code{-2**2} is \code{-4},
because the power is applied before the minus sign.

Mathematical functions live in the \code{math} module:
\begin{lstlisting}[language=Python]
import math
math.sqrt(2.0)        # 1.4142135623730951
math.exp(1.0)         # 2.718281828459045
math.log(100.0)       # natural log: 4.605170185988092
math.log10(100.0)     # 2.0
math.sin(math.pi/6)   # 0.49999999999999994  (rounding error!)
math.degrees(math.pi) # 180.0
math.hypot(3, 4)      # 5.0
\end{lstlisting}

\begin{pitfall}[Common pitfall: \texttt{\^{}} is not a power]
In \py, \code{2\^{}3} is \code{1}, not \code{8}: \code{\^{}} is the bitwise
exclusive-or operator. Use \code{2**3}. (MATLAB and Julia use \code{\^{}} for powers;
Fortran and \py use \code{**}; C, C++, Java and JavaScript need a function such as
\code{pow(2, 3)} or \code{Math.pow(2, 3)}, although JavaScript also has \code{**}.)
\end{pitfall}

\subsection*{Format output with f-strings}

Strings are written in single or double quotes. They can be joined with \code{+},
repeated with \code{*}, indexed and \emph{sliced}:
\begin{lstlisting}[language=Python]
material = "stainless steel"
len(material)          # 15
material[0]            # 's'      (first character; index 0)
material[-1]           # 'l'      (last character)
material[0:9]          # 'stainless'
material.upper()       # 'STAINLESS STEEL'
material.split()       # ['stainless', 'steel']
"steel" in material    # True
\end{lstlisting}
Strings are \emph{immutable}: methods such as \code{upper()} return a new string and leave
the original unchanged.

For output, use \term{f-strings} (formatted string literals). Put an \code{f} before the
quote and write expressions in braces. After a colon you can give a \term{format
specification}:
\begin{lstlisting}[language=Python]
v = 56.02856414093976
print(f"Terminal velocity: {v} m/s")       # all digits
print(f"Terminal velocity: {v:.2f} m/s")   # 2 decimals: 56.03
print(f"{v:10.3f}|")                       # width 10:      56.029|
print(f"{v:.3e}")                          # scientific: 5.603e+01
print(f"{1234567:,}")                      # 1,234,567
print(f"{0.0712:.1%}")                     # 7.1%
print(f"{v=:.1f}")                         # v=56.0  (handy for debugging)
\end{lstlisting}
The format mini-language (\code{10.3f}, \code{.3e}) is almost the same as the
\code{printf} format of C, which you will see again in C, C++, MATLAB, R, Julia and
Fortran-like forms.

\subsection*{Keyboard input always arrives as text}

\code{input()} prints a prompt and returns what the user types, \emph{as a string}. Convert
it before calculating:
\begin{lstlisting}[language=Python]
d = float(input("Pipe diameter in mm: "))
area = 3.141592653589793 * (d / 1000)**2 / 4
print(f"Area = {area:.3e} m^2")
\end{lstlisting}
In engineering programs, keyboard input is rare. Parameters usually come from a file or
from function arguments, which makes the program repeatable.

\section{Making decisions: \texttt{if}}

The \code{if} statement runs a block of code only when a condition is true. In \py the
block is marked by \term{indentation} (four spaces by convention) after a colon:
\begin{lstlisting}[language=Python]
Re = 45000.0
if Re < 2300:
    regime = "laminar"
elif Re < 4000:
    regime = "transitional"
else:
    regime = "turbulent"
print(f"Re = {Re:.0f}: {regime} flow")
\end{lstlisting}
Conditions use the comparison operators \code{== != < <= > >=} and can be combined with
\code{and}, \code{or} and \code{not}. \py also allows chained comparisons such as
\code{0 <= x < 1}.

\begin{keyidea}
Indentation is part of \py's grammar. In C, Java or JavaScript the braces
\code{\{ \}} define a block and indentation is only for the reader; in \py the
indentation \emph{is} the block. Never mix tabs and spaces; let your editor insert four
spaces.
\end{keyidea}

\section{Repeating: loops}

\subsection*{A \texttt{for} loop visits each item of a sequence}

A \code{for} loop takes each item of a sequence in turn:
\begin{lstlisting}[language=Python]
for material in ["steel", "aluminium", "titanium"]:
    print(material)
\end{lstlisting}
To loop over numbers, use \code{range(start, stop, step)}. The \code{stop} value is
\emph{not} included:
\begin{lstlisting}[language=Python]
# Temperature conversion table from -20 to 100 C in steps of 20
print("   C       F       K")
for T_C in range(-20, 101, 20):
    T_F = 9 / 5 * T_C + 32
    T_K = T_C + 273.15
    print(f"{T_C:4d} {T_F:7.1f} {T_K:7.2f}")
\end{lstlisting}
\begin{lstlisting}[style=output]
   C       F       K
 -20    -4.0  253.15
   0    32.0  273.15
  20    68.0  293.15
  40   104.0  313.15
  60   140.0  333.15
  80   176.0  353.15
 100   212.0  373.15
\end{lstlisting}
Two helper functions make loops cleaner: \code{enumerate} gives you the index and the item,
and \code{zip} walks through several sequences in step:
\begin{lstlisting}[language=Python]
names = ["T1", "T2", "T3"]
temps = [293.1, 305.6, 299.4]
for i, (name, T) in enumerate(zip(names, temps)):
    print(i, name, T)
\end{lstlisting}

\subsection*{A \texttt{while} loop repeats until a condition fails}

A \code{while} loop repeats as long as a condition is true. Use it when you do not know in
advance how many repetitions you need, for example in an iterative solver:
\begin{lstlisting}[language=Python]
# How many years until an investment doubles at 7% interest?
value, years = 1.0, 0
while value < 2.0:
    value *= 1.07
    years += 1
print(years)        # 11
\end{lstlisting}
\code{break} leaves a loop immediately; \code{continue} skips to the next repetition.
Every \code{while} loop needs a condition that eventually becomes false --- or a
\code{break} --- otherwise it runs forever. In an iterative method, always add a maximum
number of iterations.

\subsection*{Build up a result one step at a time}

Many calculations add up contributions one at a time. Start with a neutral value and
update it in the loop:
\begin{lstlisting}[language=Python]
# Approximate pi with the Leibniz series 4(1 - 1/3 + 1/5 - ...)
total = 0.0
for k in range(100_000):
    total += (-1)**k / (2*k + 1)
print(4 * total)    # 3.1415826535897198
\end{lstlisting}

\section{Collections}

\subsection*{A list is an ordered collection that can change}

A \term{list} is an ordered, changeable collection of values in square brackets:
\begin{lstlisting}[language=Python]
loads = [12.5, 14.0, 13.2, 15.8]      # kN
loads.append(11.9)                    # add at the end
loads[0] = 12.6                       # change an element
len(loads)                            # 5
max(loads), min(loads), sum(loads)    # (15.8, 11.9, 67.5)
sorted(loads)                         # a new sorted list
\end{lstlisting}
\term[slicing]{Slicing} extracts parts of a list with \code{[start:stop:step]}; as with
\code{range}, \code{stop} is excluded. Negative indices count from the end.
\begin{lstlisting}[language=Python]
x = [0, 10, 20, 30, 40, 50]
x[1:4]      # [10, 20, 30]
x[:3]       # [0, 10, 20]
x[3:]       # [30, 40, 50]
x[::2]      # [0, 20, 40]
x[::-1]     # [50, 40, 30, 20, 10, 0]
\end{lstlisting}

\subsection*{A comprehension builds a list in one line}

A \term{list comprehension} builds a new list from an old one in a single readable line.
It reads almost like set notation in mathematics,
$\{\, x^2 : x \in X,\ x > 0\,\}$:
\begin{lstlisting}[language=Python]
X = [-2, -1, 0, 1, 2, 3]
squares = [x**2 for x in X if x > 0]           # [1, 4, 9]
T_K = [T + 273.15 for T in [20, 25, 30]]       # [293.15, 298.15, 303.15]
\end{lstlisting}

\subsection*{Tuples are fixed, dictionaries look things up, sets remove duplicates}

A \term{tuple} is like a list but cannot be changed after creation. Tuples are used for
fixed groups of values such as coordinates, and for returning several values from a
function:
\begin{lstlisting}[language=Python]
point = (1.5, -0.3, 2.0)
x, y, z = point                 # "unpacking"
\end{lstlisting}
A \term{dictionary} maps \emph{keys} to \emph{values}. It is ideal for looking up
properties by name:
\begin{lstlisting}[language=Python]
density = {"steel": 7850, "aluminium": 2700, "water": 998}   # kg/m^3
density["steel"]                     # 7850
density["copper"] = 8960             # add an entry
for name, rho in density.items():
    print(f"{name:10s} {rho:6d} kg/m^3")
"titanium" in density                # False
density.get("titanium", 4500)        # 4500 (default if missing)
\end{lstlisting}
A \term{set} is an unordered collection without duplicates, useful for membership tests and
for removing repeated values: \code{set([1, 2, 2, 3])} is \code{\{1, 2, 3\}}.

\section{Functions}

\subsection*{Define a function once, with its units documented}

A \term{function} packages a calculation under a name so that it can be reused and
tested. It is defined with \code{def}, takes \term{parameters}, and returns a result with
\code{return}:
\begin{lstlisting}[language=Python]
import math

def terminal_velocity(m, c, g=9.81):
    """Terminal velocity (m/s) of a falling body with quadratic drag.

    m : mass in kg
    c : drag coefficient in kg/m
    g : gravitational acceleration in m/s^2 (default 9.81)
    """
    return math.sqrt(m * g / c)

print(terminal_velocity(80.0, 0.25))          # 56.02856414093976
print(terminal_velocity(80.0, 0.25, g=1.62))  # on the Moon (if it had air)
\end{lstlisting}
The string in triple quotes directly after the \code{def} line is a \term{docstring}. It
documents the function and is shown by \code{help(terminal\_velocity)} and by IDEs. Always
state the units of physical inputs and outputs.

Parameters can have \term{default values} (\code{g=9.81}), and arguments can be passed by
position or by name (\term{keyword arguments}). Keyword arguments make calls with many
parameters much easier to read.

\subsection*{A function can return several values}

A function can return a tuple, which the caller unpacks:
\begin{lstlisting}[language=Python]
def stats(values):
    n = len(values)
    mean = sum(values) / n
    var = sum((v - mean)**2 for v in values) / (n - 1)
    return mean, var**0.5

mean, std = stats([276.4, 281.9, 268.7, 273.5])
\end{lstlisting}

\subsection*{Variables inside a function stay inside}

Variables created inside a function are \term{local}: they exist only while the function
runs and are invisible outside it. A function can \emph{read} variables defined at the top
level of the file (\term{global} variables), but relying on this makes code hard to
understand and test. Pass everything a function needs as arguments, and return everything
it produces.

\subsection*{A function can be passed to another function}

In \py a function is an object like any other. You can store it in a variable or pass it
to another function. This is how general-purpose tools such as root finders and ODE
solvers work: you hand them \emph{your} function. Short one-line functions can be written
with \code{lambda}:
\begin{lstlisting}[language=Python]
def derivative(f, x, h=1e-6):
    """Central-difference approximation of f'(x)."""
    return (f(x + h) - f(x - h)) / (2 * h)

derivative(math.sin, 0.0)              # 0.9999999999998334  (cos 0 = 1)
derivative(lambda x: x**3, 2.0)        # 12.000000000789157  (exact: 12)
\end{lstlisting}

\subsection*{Newton's method shows functions and loops working together}

Newton's method finds a root of $f(x) = 0$ by repeatedly following the tangent line:
\[
x_{k+1} = x_k - \frac{f(x_k)}{f'(x_k)}.
\]
We test it on the cubic $f(x) = x^3 - 2x - 5$, which Newton himself used as an example. The
program uses functions, a loop with a stopping criterion, and an exception if the method
fails.
\codefile{Python}{python/newton.py}{Newton's method}
\outfile{python/newton.out}
Notice how quickly the step size shrinks: $10^{-1}$, $10^{-3}$, $10^{-5}$, $10^{-10}$.
The number of correct digits roughly doubles at each iteration. This is \term{quadratic
convergence}, and it is why Newton's method (in many variants) is at the heart of almost
every nonlinear solver in engineering software. The same program appears in every language
in Chapter~\ref{ch:rosetta}.

\section{Modules and packages}

A \term{module} is a \file{.py} file containing functions, classes and variables that
other programs can \code{import}. A \term{package} is a directory of modules. There are
several ways to import:
\begin{lstlisting}[language=Python]
import math                       # use as math.sqrt
import numpy as np                # standard abbreviation
from math import sqrt, pi         # use as sqrt, pi
from scipy.integrate import solve_ivp
\end{lstlisting}
Avoid \code{from module import *}: it fills your program with names whose origin is
unclear.

Your own files are modules too. If \file{fluids.py} defines \code{reynolds()}, another file
in the same folder can write \code{from fluids import reynolds}. Code at the top level of a
module runs when the module is imported. To have code that runs only when the file is
executed directly, use the idiom
\begin{lstlisting}[language=Python]
if __name__ == "__main__":
    main()
\end{lstlisting}

\subsection*{Most of what you need is already written}

The \term{standard library} is installed with \py. Useful modules for engineers include
\code{math}, \code{cmath} (complex maths), \code{statistics}, \code{random}, \code{csv},
\code{json}, \code{pathlib} (file paths), \code{datetime}, \code{time}, \code{itertools},
\code{functools}, \code{fractions}, \code{decimal} and \code{unittest}.

Everything else comes from the \term{Python Package Index} (PyPI), which hosts more than
half a million packages, installed with
\begin{lstlisting}[style=shell]
$ pip install numpy scipy matplotlib pandas
$ conda install numpy scipy matplotlib pandas     # if you use Anaconda
\end{lstlisting}
Different projects may need different versions of the same package. A \term{virtual
environment} (\code{python -m venv}, \code{conda create}, or \code{uv}) gives each project
its own isolated set of packages, and a file such as \file{requirements.txt} or
\file{environment.yml} records them so others can reproduce your setup.

\section{Errors and exceptions}

When something goes wrong at run time, \py raises an \term{exception} and prints a
\term{traceback}:
\begin{lstlisting}[language=Python]
>>> values = [1.0, 2.0]
>>> values[2]
Traceback (most recent call last):
  File "<stdin>", line 1, in <module>
IndexError: list index out of range
\end{lstlisting}
Read tracebacks from the bottom up: the last line names the error, and the lines above it
show where it happened. Common exceptions are \code{NameError} (misspelt variable),
\code{TypeError} (wrong type, e.g.\ \code{"3" + 4}), \code{ValueError} (right type,
wrong value, e.g.\ \code{float("abc")} or \code{math.sqrt(-1)}),
\code{ZeroDivisionError}, \code{IndexError}, \code{KeyError} and
\code{FileNotFoundError}.

You can catch an exception with \code{try}/\code{except} and handle it:
\begin{lstlisting}[language=Python]
text = input("Enter a temperature in K: ")
try:
    T = float(text)
except ValueError:
    print(f"'{text}' is not a number")
else:
    if T < 0:
        raise ValueError("Absolute temperature cannot be negative")
    print(f"T = {T} K")
\end{lstlisting}
\code{raise} creates an exception yourself, which is the right thing to do when a function
receives input that makes no physical sense. \code{assert} checks a condition that must
be true if the program is correct, and is useful during development:
\code{assert len(t) == len(v), "time and speed arrays differ"}.

\section{Files}

\subsection*{Open files with \texttt{with}, so that they are always closed}

Open files with a \code{with} statement, which closes them automatically even if an error
occurs:
\begin{lstlisting}[language=Python]
# Write
with open("results.txt", "w") as f:
    f.write("t v\n")
    for t, v in [(0, 0.0), (1, 9.7)]:
        f.write(f"{t} {v}\n")

# Read line by line
with open("results.txt") as f:
    header = f.readline()
    for line in f:
        t, v = line.split()
        print(float(t), float(v))
\end{lstlisting}
The mode \code{"w"} overwrites a file, \code{"a"} appends to it and \code{"r"} (the
default) reads it. The \code{pathlib} module helps with file names and folders in a way
that works on both Windows and Linux.

\subsection*{CSV is the common format for tables}

Comma-separated values (CSV) is the most common format for exchanging tables between
spreadsheets, instruments and programs. Here is the start of a file of tensile test
results. The numbers are invented for teaching; they are not measurements:
\lstinputlisting[style=output,title={\footnotesize\sffamily\color{muted}tensile.csv (first lines)},lastline=5]{code/data/tensile.csv}
The \code{csv} module reads each row as a dictionary keyed by the column names:
\codefile{Python}{python/read_csv.py}{statistics from a CSV file}
\outfile{python/read_csv.out}
For larger tables, the pandas library (Section~\ref{sec:pandas}) is more convenient.

\section{Classes and objects}

A \term{class} defines a new type that bundles data (\term{attributes}) with functions that
operate on it (\term{methods}). The special method \code{\_\_init\_\_} initialises a new
object; the first parameter of every method, by convention called \code{self}, is the
object itself. The following example models pipes and computes the pressure drop with the
Darcy--Weisbach equation
\[
\Delta p = f_D\,\frac{L}{D}\,\frac{\rho v^2}{2},
\]
where the Darcy friction factor $f_D$ is $64/\mathrm{Re}$ for laminar flow and is given by
the Swamee--Jain approximation of the Colebrook equation for turbulent flow.
\codefile{Python}{python/pipe_flow.py}{a class for pipes}
\outfile{python/pipe_flow.out}
The class keeps everything about a pipe in one place, and the loop at the bottom reads
almost like English. The special method \code{\_\_repr\_\_} controls how an object is
printed. Classes really pay off when a program has many interacting things: pipes, pumps
and tanks in a network; particles in a simulation; layers in a neural network.

\py also has \term{data classes}, which save typing when a class mainly stores data:
\begin{lstlisting}[language=Python]
from dataclasses import dataclass

@dataclass
class Material:
    name: str
    density: float          # kg/m^3
    youngs_modulus: float   # Pa

steel = Material("S275", 7850.0, 210e9)
print(steel)   # Material(name='S275', density=7850.0, youngs_modulus=210000000000.0)
\end{lstlisting}

\section{NumPy: arrays for numerical computing}
\label{sec:numpy}

\py lists can hold anything, which makes them flexible but slow for numbers. The
\term{NumPy} library provides the \code{ndarray}: a block of memory holding numbers of a
single type, laid out like a C or Fortran array, with fast compiled operations on whole
arrays. NumPy (2006, by Travis Oliphant, building on the earlier Numeric and numarray
packages) is the foundation of the entire scientific \py ecosystem.

\subsection*{Create arrays from lists, ranges and grids}

\begin{lstlisting}[language=Python]
import numpy as np

a = np.array([1.0, 2.0, 3.0])          # from a list
z = np.zeros(5)                        # [0. 0. 0. 0. 0.]
o = np.ones((2, 3))                    # 2 x 3 matrix of ones
r = np.arange(0.0, 1.0, 0.25)          # [0.   0.25 0.5  0.75]  (stop excluded)
x = np.linspace(0.0, 1.0, 5)           # [0.   0.25 0.5  0.75 1.  ]  (stop included)
I = np.eye(3)                          # 3 x 3 identity
R = np.random.default_rng(0).random(4) # 4 random numbers in [0, 1)
a.shape, a.dtype, o.ndim               # ((3,), dtype('float64'), 2)
\end{lstlisting}

\subsection*{Arithmetic applies to every element at once}

Arithmetic operators and mathematical functions work \emph{element by element}:
\begin{lstlisting}[language=Python]
t = np.linspace(0.0, 2.0, 5)
y = 5.0 * np.exp(-t) * np.cos(2 * np.pi * t)    # one line, no loop
T_C = np.array([20.0, 35.5, 50.2])
T_K = T_C + 273.15                              # [293.15 308.65 323.35]
\end{lstlisting}
Note that you must use \code{np.exp}, \code{np.sin} and so on; \code{math.exp} works
only on single numbers.

When arrays of different shapes are combined, NumPy \term{broadcasts} the smaller one.
Adding a scalar to an array is the simplest case. A more powerful example: a column of
$3$ values combined with a row of $4$ values gives a $3\times4$ table:
\begin{lstlisting}[language=Python]
D = np.array([[0.05], [0.08], [0.10]])     # 3 x 1 column: diameters (m)
Q = np.array([0.001, 0.002, 0.004, 0.008]) # 1 x 4 row:   flow rates (m^3/s)
v = Q / (np.pi * D**2 / 4)                 # 3 x 4 table of velocities
\end{lstlisting}

\subsection*{Select elements by position or by condition}

NumPy arrays are indexed from 0 and sliced like lists. Multidimensional arrays take one
index per dimension, separated by commas:
\begin{lstlisting}[language=Python]
A = np.arange(12).reshape(3, 4)   # [[0 1 2 3] [4 5 6 7] [8 9 10 11]]
A[1, 2]          # 6
A[0, :]          # first row: [0 1 2 3]
A[:, 1]          # second column: [1 5 9]
A[1:, 2:]        # [[6 7] [10 11]]
\end{lstlisting}
A comparison produces an array of booleans, which can be used to select elements. This is
called \term{boolean indexing} or masking:
\begin{lstlisting}[language=Python]
strength = np.array([276.4, 281.9, 268.7, 273.5, 279.2])
ok = strength >= 275.0              # [ True  True False False  True]
strength[ok]                        # [276.4 281.9 279.2]
np.count_nonzero(ok)                # 3
\end{lstlisting}

\begin{pitfall}[Common pitfall: slices are views]
A slice of a NumPy array is a \emph{view} into the same memory, not a copy. Changing
the slice changes the original: after \code{b = a[0:2]; b[0] = 99}, \code{a[0]} is also 99.
Use \code{a[0:2].copy()} if you need an independent array. (Slices of \py lists, in
contrast, are copies.)
\end{pitfall}

\subsection*{Summaries and linear algebra are single calls}

\begin{lstlisting}[language=Python]
x = np.array([3.0, 1.0, 4.0, 1.0, 5.0])
x.sum(), x.mean(), x.std(ddof=1), x.max(), x.argmax()   # ddof=1: sample std

A = np.array([[4.0, -2.0, 1.0],
              [-2.0, 4.0, -2.0],
              [1.0, -2.0, 4.0]])
b = np.array([11.0, -16.0, 17.0])
x = np.linalg.solve(A, b)        # solve A x = b  ->  [ 1. -2.  3.]
A @ x                            # matrix-vector product: [ 11. -16.  17.]
np.linalg.eigvalsh(A)            # eigenvalues of a symmetric matrix
np.linalg.det(A), np.linalg.norm(b)
\end{lstlisting}
The operator \code{@} is matrix multiplication; \code{*} multiplies element by element.
(In MATLAB it is the other way round: \code{*} is matrix multiplication and \code{.*} is
element-wise.)

\begin{pitfall}[Common pitfall: never invert a matrix to solve a system]
To solve $A\mathbf{x} = \mathbf{b}$, use \code{np.linalg.solve(A, b)}, not
\code{np.linalg.inv(A) @ b}. Solving directly is faster and more accurate. The same
advice holds in MATLAB (\code{A\textbackslash b}, not \code{inv(A)*b}), R, Julia and
every other language.
\end{pitfall}

\subsection*{Vectorised code is many times faster than a loop}
\label{sec:numpy-speed}

The following program evaluates the polynomial $3x^2 + 2x + 1$ at one million points, first
with a \py loop and then with one NumPy expression:
\codefile{Python}{python/numpy_speed.py}{loop versus vectorised code}
\outfile[title={\footnotesize\sffamily\color{muted}Output (times depend on your computer)}]{python/numpy_speed.out}
The vectorised version is typically 20 to 100 times faster, because the loop over the
million elements runs in compiled C code instead of the \py interpreter. This is the most
important performance rule for \py, MATLAB and R: \emph{let the library do the loops}.
When an algorithm cannot be vectorised (for example a time-stepping loop where each step
depends on the previous one), you can use a compiler such as \textbf{Numba}, which
translates a decorated \py function into machine code, or write the inner loop in C,
Fortran or Julia.

\section{Plotting with Matplotlib}

\textbf{Matplotlib} (2003, by John D.\ Hunter, a neurobiologist who wanted a free
alternative to MATLAB plotting) is the standard plotting library. The recommended way to
use it is to create a figure and axes object explicitly and call methods on the axes. The
following program produced Figure~\ref{fig:dragresult} in the guide at the front of these notes:
\codefile{Python}{python/plot_drag.py}{plotting the drag problem}
A few tips for figures in engineering reports:
\begin{itemize}
\item always label the axes \emph{with units};
\item use a legend or direct labels to identify each curve, and use different marker
  shapes as well as colours, so the figure also works in black and white;
\item show measured data as markers and model predictions as lines;
\item save as PDF or SVG (vector graphics) for reports, PNG at 300 dpi or more for slides;
\item make the font size in the figure similar to the text it sits in.
\end{itemize}

\section{SciPy: numerical methods}

\textbf{SciPy} builds on NumPy and provides well-tested numerical methods: root finding and
optimisation (\code{scipy.optimize}), integration and ODE solvers
(\code{scipy.integrate}), interpolation, linear algebra, sparse matrices, signal
processing, statistics, special functions and physical constants
(\code{scipy.constants}). Most of these routines wrap mature Fortran and C libraries.
\codefile{Python}{python/scipy_demo.py}{root finding, integration and ODEs with SciPy}
\outfile{python/scipy_demo.out}
Compare the last two lines with the Euler results in the drag program: the adaptive
Runge--Kutta method of \code{solve\_ivp} agrees with the exact solution to all printed
digits. For real work, use library solvers. Write your own only to learn how they work ---
which is exactly what we do in ENG 101.

\section{pandas: tables of data}
\label{sec:pandas}

\textbf{pandas} (2008, by Wes McKinney, then working in finance) provides the
\code{DataFrame}: a table with named columns of possibly different types, like a
spreadsheet or a database table, with powerful tools for reading, cleaning, grouping,
merging and summarising data. It was inspired by R's data frames (Chapter~\ref{ch:r}).
\codefile{Python}{python/pandas_demo.py}{a first look at pandas}
\outfile{python/pandas_demo.out}
In three lines we read the file, added two calculated columns and produced a statistical
summary per material. pandas reads and writes CSV, Excel, JSON, SQL databases, HDF5 and
Parquet files. Newer libraries such as Polars offer similar functionality with higher
speed for very large tables.

\section{Case study: the falling skydiver}

Here is the running example of these notes (page~\pageref{sec:runningexample}) in
plain \py, using lists and the \code{math} module only. Compare it line by line with the
versions in the other chapters.
\codefile{Python}{python/drag.py}{Euler solution of the drag equation}
\outfile{python/drag.out}
The error of the Euler method is largest during the acceleration phase (about
$0.9\unit{m/s}$ at $t = 6\unit{s}$) and decreases as the speed approaches the terminal
value. Halving $\Delta t$ roughly halves the error, because Euler's method is
\term{first-order accurate}. Try it.

\section{Case study: a tiny molecular dynamics simulation}

ENG 101 uses a molecular dynamics (MD) simulation of Lennard-Jones particles as a case
study. The following program shows the core of such a simulation in about 50 lines of
NumPy. Twenty-five particles interact through the Lennard-Jones potential (in reduced units)
\[
U(r) = 4\left[\left(\frac{1}{r}\right)^{12} - \left(\frac{1}{r}\right)^{6}\right],
\]
and move according to Newton's equations, integrated with the \term{velocity Verlet}
algorithm:
\begin{align*}
\mathbf{v}\!\left(t+\tfrac{\Delta t}{2}\right) &= \mathbf{v}(t) + \tfrac{\Delta t}{2}\,\mathbf{F}(t),\\
\mathbf{r}(t+\Delta t) &= \mathbf{r}(t) + \Delta t\, \mathbf{v}\!\left(t+\tfrac{\Delta t}{2}\right),\\
\mathbf{v}(t+\Delta t) &= \mathbf{v}\!\left(t+\tfrac{\Delta t}{2}\right) + \tfrac{\Delta t}{2}\,\mathbf{F}(t+\Delta t),
\end{align*}
where the mass is 1 in reduced units, so the force equals the acceleration.
\codefile{Python}{python/lj_md.py}{a minimal Lennard-Jones MD code}
\outfile{python/lj_md.out}
Kinetic and potential energy exchange continuously, but their sum, the total energy, stays
nearly constant. Checking energy conservation is the first test of any MD code. 

The line \code{d = pos[:, None, :] - pos[None, :, :]} is the hardest in the program, so
it is worth reading slowly. \code{pos} has shape $(N, 2)$: one row of $(x, y)$ per
particle. Inserting \code{None} adds an axis of length 1, so \code{pos[:, None, :]} has
shape $(N, 1, 2)$ and \code{pos[None, :, :]} has shape $(1, N, 2)$. When the two are
subtracted, NumPy broadcasts both to shape $(N, N, 2)$, and the element \code{d[i, j]} is
the vector $\mathbf{r}_i - \mathbf{r}_j$ from particle $j$ to particle $i$. One line
therefore produces all $N^2$ separation vectors without a loop. Write out the case $N = 2$
by hand if this is not yet clear. Computing all pairs at once is fast for small $N$; production codes such as LAMMPS (written in C++) use neighbour lists so that
the cost grows only like $N$, and they run on thousands of processors.

\section{Recursion}

A function may call itself. This is called \term{recursion}, and it is natural for
problems that are defined in terms of smaller versions of themselves. The factorial
$n! = n \cdot (n-1)!$ with $0! = 1$ is the classic example:
\begin{lstlisting}[language=Python]
def factorial(n):
    if n == 0:                 # base case: stops the recursion
        return 1
    return n * factorial(n - 1)

factorial(20)                  # 2432902008176640000
\end{lstlisting}
Every recursive function needs a \emph{base case} that does not call the function again;
without it, the calls never end. \py limits the depth of recursion (by default to about
1000 calls) and raises a \code{RecursionError} when the limit is exceeded. In engineering
codes, recursion appears in adaptive integration (split an interval in two if the error is
too large), in tree structures such as the octrees used in particle simulations and mesh
generation, and in divide-and-conquer algorithms such as the fast Fourier transform. For
simple repetitions a loop is usually clearer and faster: the naive recursive Fibonacci
function \code{fib(n) = fib(n-1) + fib(n-2)} recomputes the same values so often that
\code{fib(30)} makes more than a million function calls, whereas a loop needs thirty steps.

\section{Debugging}

Every programmer spends a large part of their time finding errors. Debugging is a skill,
and it can be learnt. A systematic approach:
\begin{enumerate}
\item \textbf{Reproduce the error.} Find the smallest input that shows the problem. A bug
  that appears only sometimes is often caused by uninitialised data, random numbers without
  a fixed seed, or the order in which notebook cells were run.
\item \textbf{Read the message.} A traceback names the type of error, the file and the line.
  Look at that line, and at the line before it.
\item \textbf{Form a hypothesis and test it.} ``The velocity becomes negative because the
  time step is too large.'' Then check: print the value, or change the time step.
\item \textbf{Divide and conquer.} If a long calculation gives a wrong result, check
  intermediate values halfway through. Is the error before or after that point? Repeat.
\item \textbf{Explain the code to someone else} --- a classmate, or even a rubber duck on
  your desk. Explaining forces you to look at what the code actually does, not what you
  meant it to do. (Programmers really call this ``rubber duck debugging''.)
\end{enumerate}

\subsection*{Print first, then reach for the debugger}

The simplest tool is a well-placed \code{print}. The f-string form \code{print(f"\{v=\}")}
prints both the name and the value of a variable. For more than a few prints, use a
\term{debugger}. In VS Code, Spyder or PyCharm, click to the left of a line number to set a
\term{breakpoint}, then start the program in debug mode. Execution stops at the breakpoint,
and you can inspect every variable, step through the code line by line (\emph{step over}),
enter a function call (\emph{step into}), and continue to the next breakpoint. From the
command line, insert the built-in function \code{breakpoint()} where you want the program to
stop; this opens \py's text-based debugger \code{pdb}:
\begin{lstlisting}[language=Python]
for n in range(n_steps):
    if v[n] < 0:
        breakpoint()        # stops here; type p v[n] to print, c to continue, q to quit
    v[n + 1] = v[n] + dt * acceleration(v[n])
\end{lstlisting}
In Jupyter, the magic command \code{\%debug} run after an exception opens the debugger at
the place where the error occurred.

\begin{tryit}
Introduce a deliberate bug into \file{drag.py}: change \code{v[n + 1]} to \code{v[n]} in the
update line. Run the program. What do you see? Use the debugger to step through the first
three iterations and explain the output.
\end{tryit}

\section{Working with notebooks}

Jupyter notebooks are excellent for exploration and for explaining an analysis, and ENG 101
uses them for case studies. A few habits make them much more reliable:
\begin{itemize}
\item Put all imports and parameters in the first cells, so that a reader sees at once
  what the notebook depends on.
\item Run cells from top to bottom. Before saving or sharing, use ``Restart kernel and run
  all'' to check that the notebook works in order.
\item Keep cells short; move long functions into a \file{.py} module and import them. Code
  in modules can be tested and reused by other notebooks.
\item Use Markdown cells to say \emph{why} you do something, and \LaTeX\ notation
  (\code{\$v\_t = \textbackslash sqrt\{mg/c\}\$}) to write equations.
\item Keep raw data files unchanged; save processed data and figures under new names.
\end{itemize}

\section{Writing good Python}

\subsection*{Follow the common style so that others can read your code}

\py has an official style guide, \textbf{PEP~8}. Its main points: four spaces per
indentation level; lines shorter than about 80--100 characters; \code{lower\_case} names for
variables and functions, \code{CapitalisedWords} for classes and \code{UPPER\_CASE} for
constants; spaces around operators; and imports at the top of the file. Tools such as
\code{ruff} and \code{black} check and format code automatically.

\subsection*{Type hints document what a function expects}

Since \py~3.5 you can annotate the expected types of arguments and return values:
\begin{lstlisting}[language=Python]
import numpy as np

def kinetic_energy(m: float, v: np.ndarray) -> float:
    return 0.5 * m * float(np.sum(v**2))
\end{lstlisting}
The interpreter ignores these hints at run time, but editors use them to warn you about
mistakes before you run the code, and tools such as \code{mypy} check them systematically.
They are also excellent documentation.

\subsection*{A test compares a result with an answer you already know}

A program you have not tested is a program you should not trust. The simplest test
compares a computed result with a known answer:
\begin{lstlisting}[language=Python]
# test_drag.py -- run with:  pytest
import math
from drag_lib import terminal_velocity

def test_terminal_velocity():
    assert math.isclose(terminal_velocity(80.0, 0.25), 56.0285641, rel_tol=1e-7)

def test_terminal_velocity_scales_with_sqrt_mass():
    v1 = terminal_velocity(80.0, 0.25)
    v4 = terminal_velocity(320.0, 0.25)
    assert math.isclose(v4, 2 * v1)
\end{lstlisting}
The \code{pytest} tool finds functions whose names start with \code{test\_} and runs them.
Good tests use exact solutions, limiting cases, conservation laws (as in the MD example)
and \emph{scaling} relations like the second test above.

\section{Strengths and weaknesses}
\label{sec:python-strengths}

\begin{center}
\small
\begin{tabularx}{\linewidth}{@{}>{\raggedright\arraybackslash}X>{\raggedright\arraybackslash}X@{}}
\toprule
\textbf{Strengths} & \textbf{Weaknesses}\\
\midrule
Easy to learn and read; short programs & Pure-\py loops are slow (10--100$\times$ slower than C)\\
Free, open source, runs everywhere & Dynamic typing hides some errors until run time\\
Huge ecosystem for science, data and AI & Package and environment management can be confusing\\
Interactive use and notebooks & Not suitable for hard real-time or small embedded systems\\
Excellent glue to C, C++, Fortran and to engineering software & The Global Interpreter Lock limits multi-threaded speed (being removed in recent versions)\\
Very large community; most AI tools know it well & Deployment as a standalone application is clumsy\\
\bottomrule
\end{tabularx}
\end{center}

\begin{engnote}[Engineering note: Python is often the conductor, not the orchestra]
In a typical research or industrial workflow, \py prepares input files, launches a
simulation written in Fortran or C++, reads the results, analyses them with NumPy and
pandas, and makes the plots for the report. \py coordinates; the compiled programs do the
heavy computing. Knowing a little about those compiled languages (the next part of these
notes) makes you much more effective in that role.
\end{engnote}

\section{Further reading}
\begin{itemize}
\item The official \py tutorial, \url{https://docs.python.org/3/tutorial/}.
\item A.~B.~Downey, \emph{Think Python}, 3rd ed., O'Reilly, 2024 (free online).
\item J.~VanderPlas, \emph{Python Data Science Handbook}, 2nd ed., O'Reilly, 2022 (free online).
\item J.~Kiusalaas, \emph{Numerical Methods in Engineering with Python 3}, Cambridge
  University Press, 2013.
\item The NumPy, SciPy, Matplotlib and pandas documentation, each with a ``getting
  started'' guide.
\end{itemize}

\begin{exercises}
\item Write a script that asks for a mass and a drag coefficient and prints the terminal
  velocity with two decimals. What happens if the user types a negative drag coefficient?
  Add a check that raises a \code{ValueError} with a helpful message.
\item Write a function \code{celsius\_to\_fahrenheit(T)} and use a list comprehension to
  convert the list \code{[-40, 0, 37, 100]}. Which value is the same on both scales?
\item Modify \file{drag.py} so that it runs with $\Delta t = 0.5$, $0.25$, $0.125$ and
  $0.0625\unit{s}$ and prints the maximum absolute error for each. Show that the error is
  roughly proportional to $\Delta t$.
\item Replace the lists in \file{drag.py} with NumPy arrays and compute the exact solution
  for all time points with a single vectorised expression.
\item Use \code{np.linalg.solve} to find the currents in a circuit with three loops whose
  mesh equations you write down from a circuit diagram of your choice. Check the answer by
  substituting it back into the equations.
\item Extend the \code{Pipe} class with a method \code{head\_loss(Q)} that returns the
  pressure drop expressed as a height of water, $h = \Delta p/(\rho g)$. Then compute how the
  pressure drop of a 50~m pipe with $D = 0.05\unit{m}$ varies with flow rate for $Q$ between
  $0.0005$ and $0.01\unit{m^3/s}$, and plot it on log--log axes. What slope do you expect?
\enlargethispage{2\baselineskip}
\item Read \file{tensile.csv} with pandas, compute the yield load (in kN) of each specimen
  from its diameter, and find the specimen with the largest yield load.
\item In \file{lj\_md.py}, change the time step to $0.02$ and then $0.05$. What happens to
  the total energy? Explain why a time step that is too large is dangerous in a
  simulation.
\end{exercises}
